\documentclass[12pt,a4paper,reqno]{amsart}

% language
\usepackage[greek,english]{babel}
\usepackage[utf8]{inputenc}
\usepackage{alphabeta}

% change default names to greek
\addto\captionsenglish{
    \renewcommand{\contentsname}{Περιεχόμενα}
    \renewcommand{\refname}{Βιβλιογραφία}
    \renewcommand{\datename}{Ημερομηνία:}
    \renewcommand{\urladdrname}{Ιστοσελίδα}
}

% math 
\usepackage{amsmath,amsthm,amssymb,amscd}

% font
\usepackage[cal=euler]{mathalfa}
\usepackage{libertinus-type1}
% \usepackage{txfonts} % for upright greek letters
\usepackage{bm} % for bold symbols
\usepackage{bbm} % for the simply-looking bb symbols

% miscellaneous 
\usepackage{changepage} %for indenting environments
\usepackage{csquotes} % example: \textcquote{}
\usepackage{blkarray}
\setcounter{MaxMatrixCols}{20} % default for pmatrix is 10!!
\usepackage{ytableau}
\usepackage{array} %needed to increase the vertical length in a tabular

% drawing
\usepackage{tikz,tikz-cd}
\usetikzlibrary{shapes.misc, patterns, matrix, calc, intersections,positioning,arrows.meta}
\usepackage{graphics,graphicx}
\usepackage{float} % provides enhanced control and customization options for floating objects such as figures and tables

% colors
\usepackage{xcolor}
\definecolor{darkcandyapplered}{rgb}{0.64, 0.0, 0.0}
\definecolor{midnightblue}{rgb}{0.1, 0.1, 0.44}
\definecolor{mylightblue}{HTML}{336699}
\definecolor{burntorange}{rgb}{0.8, 0.33, 0.0}
\definecolor{iceberg}{rgb}{0.44, 0.65, 0.82}
\definecolor{applegreen}{rgb}{0.55, 0.71, 0.0}
\definecolor{canaryyellow}{rgb}{1.0, 0.94, 0.0}
\definecolor{lightgray}{rgb}{0.83, 0.83, 0.83}

% hrefs
\usepackage{hyperref}
\usepackage[noabbrev,capitalize]{cleveref}
\hypersetup{
    pdftoolbar=true,        
    pdfmenubar=true,        
    pdffitwindow=false,     
    pdfstartview={FitH},  % fits the width of the page to the window
    pdftitle={},
    pdfauthor={},
    pdfsubject={},
    pdfkeywords={},
    pdfnewwindow=true,  % links in new window
    colorlinks=true,  % false: boxed links; true: colored links
    linkcolor=darkcandyapplered,   % color of internal links
    citecolor=midnightblue,  % color of links to bibliography
    urlcolor=cyan,  % color of external links
    linktocpage=true  % changes the links from the section body to the page number
    }

% geometry
\textwidth=16cm 
\textheight=21cm 
\hoffset=-55pt 
\footskip=25pt

% thm envs (you might need to change the path)
\theoremstyle{definition}
\newtheorem{theorem}{Θεώρημα}
\newtheorem{proposition}[theorem]{Πρόταση}
\newtheorem{lemma}[theorem]{Λήμμα}
\newtheorem{corollary}[theorem]{Πόρισμα}
\newtheorem{conjecture}[theorem]{Εικασία}
\newtheorem{problem}[theorem]{Πρόβλημα}
\newtheorem*{claim}{Ισχυρισμός}
\newtheorem{observation}[theorem]{Παρατήρηση}
\newtheorem{definition}[theorem]{Ορισμός}
\newtheorem{question}[theorem]{Ερώτημα}
\newtheorem*{questions}{Ερωτήματα}
\newtheorem*{que}{Ερώτημα}
\newtheorem{example}[theorem]{Παράδειγμα}
\newtheorem{exercise}{Άσκηση}
\newtheorem*{zorn}{Λήμμα του Zorn}
\newtheorem*{example_6.6}{Παράδειγμα~6.6}

\theoremstyle{remark}
\newtheorem*{remark}{Παρατήρηση}

% fixes the correct numbering of environments
\numberwithin{theorem}{section}
\numberwithin{exercise}{section}
\numberwithin{equation}{section}

% math ops 
% fields
\renewcommand{\AA}{\mathbbmss{A}}
\newcommand{\NN}{\mathbbmss{N}} 
\newcommand{\ZZ}{\mathbbmss{Z}} 
\newcommand{\QQ}{\mathbbmss{Q}} 
\newcommand{\RR}{\mathbbmss{R}} 
\newcommand{\CC}{\mathbbmss{C}} 
\newcommand{\KK}{\mathbbmss{K}} 
\newcommand{\FF}{\mathbbmss{F}} 
\newcommand{\HH}{\mathbbmss{H}} 
\newcommand{\VV}{\mathbbmss{V}} 
\newcommand{\II}{\mathbbmss{I}} 

% symmetric group
\newcommand{\fS}{\mathfrak{S}}  

% calligraphic 
\newcommand{\aA}{\mathcal{A}} 
\newcommand{\bB}{\mathcal{B}}
\newcommand{\cC}{\mathcal{C}}
\newcommand{\dD}{\mathcal{D}}
\newcommand{\eE}{\mathcal{E}}
\newcommand{\fF}{\mathcal{F}}
\newcommand{\hH}{\mathcal{H}}
\newcommand{\iI}{\mathcal{I}}
\newcommand{\lL}{\mathcal{L}}
\newcommand{\oO}{\mathcal{O}}
\newcommand{\pP}{\mathcal{P}}
\newcommand{\qQ}{\mathcal{Q}}
\newcommand{\sS}{\mathcal{S}}
\newcommand{\mM}{\mathcal{M}}
\newcommand{\uU}{\mathcal{U}}
\newcommand{\zZ}{\mathcal{Z}}
\newcommand{\vV}{\mathcal{V}}

% bold
\newcommand{\bfa}{\pmb{a}}
\newcommand{\bfe}{\mathbf{e}}
\newcommand{\bfF}{\pmb{F}}
\newcommand{\bfR}{\pmb{R}}
\newcommand{\bfv}{\mathbf{v}}
%\newcommand{\bfx}{\bm{x}}
%\newcommand{\bfx}{\mathbf{x}} 
\newcommand{\bfx}{\pmb{x}}
\newcommand{\bfX}{\pmb{X}}
\newcommand{\bfy}{\pmb{y}}
\newcommand{\bfz}{\pmb{z}}

% roman
\newcommand{\rmA}{\mathrm{A}}
\newcommand{\rmB}{\mathrm{B}}
\newcommand{\rmC}{\mathrm{C}}
\newcommand{\rmD}{\mathrm{D}} 
\newcommand{\rmF}{\mathrm{F}} 
\newcommand{\rmH}{\mathrm{H}} 
\newcommand{\rmh}{\mathrm{h}} 
\newcommand{\rmI}{\mathrm{I}} 
\newcommand{\rmK}{\mathrm{K}}
\newcommand{\rmM}{\mathrm{M}}
\newcommand{\rmP}{\mathrm{P}}  
\newcommand{\rmp}{\mathrm{p}}  
\newcommand{\rmQ}{\mathrm{Q}}  
\newcommand{\rmR}{\mathrm{R}}
\newcommand{\rmS}{\mathrm{S}}
\newcommand{\rmT}{\mathrm{T}}
\newcommand{\rmU}{\mathrm{U}}
\newcommand{\rmV}{\mathrm{V}}
\newcommand{\rmY}{\mathrm{Y}}
\newcommand{\rmZ}{\mathrm{Z}}
\newcommand{\rmz}{\mathrm{z}}

% arrows and symbols 
\renewcommand{\to}{\rightarrow}
\newcommand{\toto}{\longrightarrow}
\newcommand{\mapstoto}{\longmapsto}
\newcommand{\then}{\Rightarrow}
\newcommand{\IFF}{\Leftrightarrow}
\newcommand{\tl}{\tilde}
\newcommand{\wtl}{\widetilde}
\newcommand{\ol}{\overline}
\newcommand{\ul}{\underline}
\newcommand{\oldemptyset}{\emptyset}
\renewcommand{\emptyset}{\varnothing}
\DeclareMathSymbol{\Arg}{\mathbin}{AMSa}{"39} % for arguments 
\newcommand{\onto}{\ensuremath{\twoheadrightarrow}}
\newcommand{\tle}{\trianglelefteq}
\newcommand{\tge}{\trianglerighteq}

% custom ops
\newcommand{\rmKer}{\mathrm{Ker}}
\newcommand{\rmIm}{\mathrm{Im}}
\newcommand{\lcm}{\mathrm{lcm}}
\newcommand{\GL}{\mathrm{GL}}
\newcommand{\ev}{\mathrm{ev}}
\newcommand{\End}{\mathrm{End}}
\newcommand{\rmid}{\mathrm{id}}
\newcommand{\rmspan}{\mathrm{span}}
\newcommand{\Hom}{\mathrm{Hom}}
\newcommand{\Ann}{\mathrm{Ann}}
\newcommand{\Set}{\pmb{\mathrm{Set}}}
\newcommand{\Rmod}{\pmb{\mathrm{mod}}}
\newcommand{\rmrank}{\mathrm{rank}}
\newcommand{\mSpec}{\mathrm{mSpec}}
\newcommand{\Spec}{\mathrm{Spec}}
\newcommand{\tor}{\mathrm{tor}}
\newcommand{\Pol}[1]{F[x_1, x_2, \dots, x_{#1}]}

% absolute value symbol
\usepackage{mathtools} 
\DeclarePairedDelimiter\abs{\lvert}{\rvert}%
\DeclarePairedDelimiter\norm{\lVert}{\rVert}%
\makeatletter
\let\oldabs\abs
\def\abs{\@ifstar{\oldabs}{\oldabs*}}

% tensor symbol
\newcommand{\tensor}[1]{%
  \mathbin{\mathop{\otimes}\limits_{#1}}%
}

% permutation cycle notation
\ExplSyntaxOn
\NewDocumentCommand{\cycle}{ O{\;} m }
 {
  (
  \alec_cycle:nn { #1 } { #2 }
  )
 }

\seq_new:N \l_alec_cycle_seq
\cs_new_protected:Npn \alec_cycle:nn #1 #2
 {
  \seq_set_split:Nnn \l_alec_cycle_seq { , } { #2 }
  \seq_use:Nn \l_alec_cycle_seq { #1 }
 }
\ExplSyntaxOff

% setminus symbol
\newcommand{\mysetminusD}{\hbox{\tikz{\draw[line width=0.6pt,line cap=round] (3pt,0) -- (0,6pt);}}}
\newcommand{\mysetminusT}{\mysetminusD}
\newcommand{\mysetminusS}{\hbox{\tikz{\draw[line width=0.45pt,line cap=round] (2pt,0) -- (0,4pt);}}}
\newcommand{\mysetminusSS}{\hbox{\tikz{\draw[line width=0.4pt,line cap=round] (1.5pt,0) -- (0,3pt);}}}
\newcommand{\sm}{\mathbin{\mathchoice{\mysetminusD}{\mysetminusT}{\mysetminusS}{\mysetminusSS}}}

% miscellaneous commands
\newcommand{\defn}[1]{{\color{mylightblue}{#1}}}
\newcommand{\toDo}{{\bf\color{red} TODO}}
\newcommand{\toCite}{{\bf\color{green} CITE}}
\newcommand*{\vertbar}{\rule[-1ex]{0.5pt}{2.5ex}} % for matrices with column vectors
\newcommand*{\horzbar}{\rule[.5ex]{2.5ex}{0.5pt}} % for matrices with row vectors
\newcommand{\myblue}[1]{{\color{iceberg}{#1}}}
\newcommand{\myorange}[1]{{\color{burntorange}{#1}}}
\newcommand{\mygreen}[1]{{\color{applegreen}{#1}}}
\newcommand{\myred}[1]{{\color{darkcandyapplered}{#1}}}

% ferrer's diagram
\newcommand{\fdiagram}[1]{
    \begin{tikzpicture}[scale=.7]
        \fill foreach \Z [count=\Y] in {#1}
        {foreach \X in {1,...,\Z} 
        {(\X,-\Y) circle[radius=3pt]}};
    \end{tikzpicture}
}

%
\usepackage{pgfplots}
\pgfplotsset{compat=1.18}

%
\makeatletter
\def\mathcolor#1#{\@mathcolor{#1}}
\def\@mathcolor#1#2#3{%
  \protect\leavevmode
  \begingroup
    \color#1{#2}#3%
  \endgroup
}
\makeatother

% restriction
\newcommand\restr[2]{{% we make the whole thing an ordinary symbol
  \left.\kern-\nulldelimiterspace % automatically resize the bar with \right
  #1 % the function
  \littletaller % pretend it's a little taller at normal size
  \right|_{#2} % this is the delimiter
  }}
\newcommand{\littletaller}{\mathchoice{\vphantom{\big|}}{}{}{}}

% 
\newenvironment{nouppercase}{%
  \let\uppercase\relax%
  \renewcommand{\uppercasenonmath}[1]{}}{}

% titlepage
\title{226: Θεωρία Δακτυλίων και Modules}
\author[Β.~Δ. Μουστακας]{Βασίλης Διονύσης Μουστάκας \\ Πανεπιστήμιο Κρήτης}
\date{5 Μαΐου 2026}
% \urladdr{\href{https://sites.google.com/view/vasmous}{https://sites.google.com/view/vasmous}}

\begin{document}

\begingroup
\def\uppercasenonmath#1{} % this disables uppercase title
\let\MakeUppercase\relax % this disables uppercase authors
\maketitle
\endgroup

\setcounter{section}{9}
\setcounter{theorem}{2}
\setcounter{equation}{0}
\begin{center}
    \textbf{9. Αλγεβρική Γεωμετρία: Εισαγωγή
} (Συνέχεια)
\end{center}

\begin{example}{\rm(Συνέχεια)}
    \label{ex:variety}
    \leavevmode
    \begin{itemize}
        \item[(δ)] Τα μονοσύνολο $\{\bfa\}$, για κάθε $\bfa \in \AA^n$, είναι πολλαπλότητα, διότι \[\{\bfa\} = \vV(\{x_1-a_1, x_2-a_2, \dots, x_n-a_n\}).\]
        \item[(ε)] Αν το $S$ είναι ένα σύνολο \emph{ομογενών} γραμμικών πολυωνύμων, δηλαδή της μορφής 
        \[
        f_i(x_1, x_2, \dots, x_n) = a_{i1}x_1 + a_{i2}x_2 + \cdots + a_{in}x_n
        \]
        τότε η πολλαπλότητα $\vV(S)$ είναι ένας υπόχωρος του $F^n$, ενώ αν επιτρέψουμε \emph{μη ομογενή} γραμμικά πολυώνυμα, δηλαδή πολυώνυμα της μορφής 
        \[
        f_i(x_1, x_2, \dots, x_n) = a_{i1}x_1 + a_{i2}x_2 + \cdots + a_{in}x_n - b_i
        \]
        τότε η πολλαπλότητα $\vV(S)$ είναι είτε $\emptyset$, είτε μεταφορές υπόχωρων του $F^n$. Αυτό εξηγεί τον όρο αφφινικές ή ομοπαραλληλικές. Για παράδειγμα,  οι $\vV(y-x)$ και $\vV(y-x-1)$
        \[
        \begin{tikzpicture}
        \begin{axis}[
            axis lines = middle,
            axis line style={-latex, gray!50},
            xlabel = {$x$},
            ylabel = {$y$},
            label style = {gray!50},
            xmin=-2, xmax=2,
            ymin=-2, ymax=2,
            samples=100,
            xticklabels=\empty,
            yticklabels=\empty,
            ticks=none
        ]
    
        % function
        \addplot[thick] {x}; 
        \end{axis}

        \begin{axis}[
            at={(8cm,0)},
            axis lines = middle,
            axis line style={-latex, gray!50},
            xlabel = {$x$},
            ylabel = {$y$},
            label style = {gray!50},
            xmin=-2, xmax=2,
            ymin=-2, ymax=2,
            samples=100,
            xticklabels=\empty,
            yticklabels=\empty,
            ticks=none
        ]
    
        % function
        \addplot[thick] {x-1}; 
        \end{axis}
        \end{tikzpicture}
        \]
        
        \item[(στ)] Οι πολλαπλότητες της μορφής $\vV(f) \subseteq \AA^n$, για κάποιο (μη σταθερό) $f \in F[x_1,x_2,$ $\dots, x_n]$ ονομάζονται \emph{υπερεπιφάνειες} (hypersurfaces). Ειδικότερα,
        \begin{itemize}
            \item αν $n=2$, τότε ονομάζονται \emph{αλγεβρικές καμπύλες} (algebraic curves)
            \item αν $n=3$, τότε ονομάζονται \emph{επιφάνειες} (surfaces)
            \item αν $\deg(f) = 1$, τότε ονομάζονται \emph{υπερεπίπεδα} (hyperplanes).
        \end{itemize}
        Οι αλγεβρικές καμπύλες της μορφής 
        \[
        f(x,y) = a_{20}x^2 + a_{02}y^2 + a_{11}xy + a_{10}x + a_{01}y + a_{00} 
        \]
        για παραμέτρους $a_{20}, a_{02}, a_{11}, a_{10}, a_{01}, a_{00} \in F$ ονομάζονται \emph{κωνικές τομές} (conic sections) και μελετώνται στην Αναλυτική και την Ευκλείδια Γεωμετρία. Παραδείγματα (πραγματικών) κωνικών τομών αποτελούν ο κύκλος, η παραβολή, η έλλειψη και η υπερβολή
        \[
        \begin{tikzpicture}
            % --- Plot 1: Circle (x^2 + y^2 = 1) ---
            \begin{axis}[
                axis lines = middle,
                axis line style={-latex, gray!50},
                xlabel = {$x$},
                ylabel = {$y$},
                label style={gray!50},
                xmin=-1.5, xmax=1.5,
                ymin=-1.5, ymax=1.5,
                samples=200,
                xticklabels=\empty,
                yticklabels=\empty,
                ticks=none,
                unit vector ratio=1 1 % Ensures the circle isn't an ellipse
            ]
                % Parametric plot for a smooth circle
                \addplot[thick, domain=0:360] ({cos(x)}, {sin(x)});
            \end{axis}

            % --- Plot 2: Parabola (y = x^2) ---
            \begin{axis}[
                at={(8cm,0)},
                axis lines = middle,
                axis line style={-latex, gray!50},
                xlabel = {$x$},
                ylabel = {$y$},
                label style = {gray!50},
                xmin=-1.5, xmax=1.5,
                ymin=-2, ymax=4,
                samples=100,
                xticklabels=\empty,
                yticklabels=\empty,
                ticks=none
            ]
    
        % function
        \addplot[thick] {x^2}; 
        \end{axis}
        \end{tikzpicture}
        \]

        \[
        \begin{tikzpicture}
            \begin{axis}[
            at={(8cm,0)},
            axis lines = middle,
            axis line style={-latex, gray!50},
            xlabel = {$x$},
            ylabel = {$y$},
            label style={gray!50},
            xmin=-3, xmax=3,
            ymin=-3, ymax=3,
            samples=100,
            xticklabels=\empty,
            yticklabels=\empty,
            ticks=none
            ]
            % Positive branch
            \addplot[thick, domain=0.33:2.8] {1/x};
            % Negative branch
            \addplot[thick, domain=-2.8:-0.33] {1/x};
            \end{axis}

            \begin{axis}[
            axis lines = middle,
            axis line style={-latex, gray!50},
            xlabel = {$x$},
            ylabel = {$y$},
            label style={gray!50},
            xmin=-2.5, xmax=2.5,
            ymin=-1.5, ymax=1.5,
            samples=200,
            xticklabels=\empty,
            yticklabels=\empty,
            ticks=none,
            unit vector ratio=1 1,
            title style={gray!50, yshift=-10pt}
        ]
            % Parametric plot: x = 2cos(t), y = 1sin(t)
            \addplot[thick, domain=0:360] ({2*cos(x)}, {sin(x)});
            \end{axis}  
        \end{tikzpicture}
        \]
        με πολυώνυμα $x^2 + y^2 -1, \ x^2/4 + y^2 -1, \ x^2-y$ και $xy - 1$, αντίστοιχα.
        
        Οι αλγεβρικές καμπύλες της μορφής 
        \[
        f(x,y) = y^2 - x^3 - ax - b
        \]
        για παραμέτρους $a, b \in F$ ονομάζονται \emph{ελλειπτικές καμπύλες} (elliptic curves) και η μελέτη τους έπαιξε καθοριστικό ρόλο στην απόδειξη του τελευταίου θεωρήματος του Fermat. Παραδείγματα (πραγματικών) ελλειπτικών καμπυλών αποτελούν 
        \[
        \begin{tikzpicture}
        \begin{axis}[
            axis lines = middle,
            axis line style={-latex, gray!50},
            xlabel = {$x$},
            ylabel = {$y$},
            label style={gray!50},
            % Set limits to show the curve clearly.
            % We choose xmin=-1.2 because x must be >= -1.
            xmin=-1.2, xmax=2.5,
            ymin=-4, ymax=4,
            samples=100,
            xticklabels=\empty,
            yticklabels=\empty,
            ticks=none,
            title style={gray!50} % Title color
        ]

        % The upper branch of the curve: y = sqrt(x^3 + 1)
        % We start the domain precisely at x=-1 to avoid imaginary numbers.
        \addplot[thick, domain=-1:2.3] {sqrt(x^3 + 1)};
        
        % The lower branch of the curve: y = -sqrt(x^3 + 1)
        \addplot[thick, domain=-1:2.3] {-sqrt(x^3 + 1)};
        
        \end{axis}

        \begin{axis}[
            at={(8cm,0)},
            axis lines = middle,
            axis line style={-latex, gray!50},
            xlabel = {$x$},
            ylabel = {$y$},
            label style={gray!50},
            xmin=-1.5, xmax=3,
            ymin=-4, ymax=4,
            samples=200, % Higher samples help smooth the oval edges
            xticklabels=\empty,
            yticklabels=\empty,
            ticks=none
        ]

        % --- Component 1: The Oval (-1 <= x <= 0) ---
        \addplot[thick, domain=-1:0] {sqrt(x^3 - x)};
        \addplot[thick, domain=-1:0] {-sqrt(x^3 - x)};

        % --- Component 2: The Infinite Branch (x >= 1) ---
        \addplot[thick, domain=1:2.5] {sqrt(x^3 - x)};
        \addplot[thick, domain=1:2.5] {-sqrt(x^3 - x)};

        \end{axis}
        \end{tikzpicture}
        \]
        για τις τιμές $(a,b) =(0,1)$ και $(a,b) = (-1,0)$, αντίστοιχα.

        \item[(ζ)] Διαφορετικά σύνολα πολυωνύμων μπορεί να ορίσουν τις ίδιες πολλαπλότητες. Για παράδειγμα, τα $(x)$ και $(x^2)$ είναι δύο διαφορετικά ιδεώδη του $F[x]$, αλλά 
        \[
        \vV(x) = \vV(x^2) = \{0_F\} \ \subseteq \ \AA^1.
        \]

        \item[(η)] Η αλγεβρική κλειστότητα του $F$ παίζει σημαντικό ρόλο. Για παράδειγμα, τα πολυώ\-νυμα 
        \[
        1, \quad 1 + x^2, \quad 1 + x^2 + x^4
        \]
        του $\RR[x]$ παράγουν διαφορετικά ιδεώδη, αλλά δεν έχουν ρίζες στο $\RR$ και γι αυτό 
        \[
        \vV(1) = \vV(1+x^2) = \vV(1 + x^2 + x^4) = \emptyset \subset \AA_\RR^1.
        \]
        Όμως, ως υποσύνολα του $\AA_\CC^1$ έχουμε 
        \begin{align*}
            \vV(1) &= \emptyset \\
            \vV(1 + x^2) &= \{\pm i\} \\
            \vV(1 + x^2 + x^4) &= \{\pm \frac{1}{2} \pm i\frac{\sqrt{3}}{2}\}. 
        \end{align*}
        \item[(θ)] Δεν είναι όλα τα υποσύνολα του αφφινικού χώρου πολλαπλότητες. Για παράδειγμα, το $\ZZ \subseteq \AA_\RR^1$ δεν είναι πολλαπλότητα, διότι κάθε πολυώνυμο του $\RR[x]$ έχει πεπερασμένο πλήθος ριζών (γιατί;).
    \end{itemize}
\end{example}

\begin{proposition}
    \label{prop:zero_locus}
    Έστω $\Lambda$ ένα σύνολο και $I, J$ και $I_\lambda$ ιδεώδη του $\Pol{n}$, για κάθε $\lambda \in \Lambda$. Τότε 
    \begin{align}
        \vV(I) \cup \vV(J) &= \vV(I\cap{J}) = \vV(IJ) \label{eq:zero_locus_cup} \\
        \bigcap_{\lambda \in \Lambda}\vV(I_\lambda) &= \vV\left(
            \sum_{\lambda \in \Lambda} I_\lambda
        \right). \label{eq:zero_locus_cap}
    \end{align}
\end{proposition}

\begin{proof}[Απόδειξη]
    Για την Ταυτότητα \eqref{eq:zero_locus_cup}, από την μια μεριά έχουμε $
    IJ \ \subseteq \ I\cap{J} \ \subseteq \ I, J$ (βλ. Πρόταση 1.6) και γι αυτό $\vV(I) \cup \vV(J) \subseteq \vV(I\cap{J}) \subseteq \vV(IJ)$.
    Από την άλλη, αν $\bfa \notin \vV(I) \cup \vV(J)$, τότε υπάρχουν $f \in I$ και $g \in J$ τέτοια ώστε $f(\bfa) \neq 0$ και $g(\bfa) \neq 0$. Συνεπώς, $f(\bfa)g(\bfa) \neq 0$ και γι αυτό $\bfa \notin \vV(IJ)$. Άρα, $\vV(IJ) \subseteq \vV(I\cup J)$ και η ισότητα έπεται.

    Όμοια, για την Ταυτότητα \eqref{eq:zero_locus_cap}, επειδή $I_\mu \subset \sum_{\lambda \in \Lambda} I_\lambda$, για κάθε $\mu \in \Lambda$, έπεται ότι 
    \[     
    \vV\left(
            \sum_{\lambda \in \Lambda} I_\lambda
        \right)\subseteq I_\mu,
    \] 
    για κάθε $\mu \in \Lambda$. Επομένως, 
    \[
    \vV\left(
            \sum_{\lambda \in \Lambda} I_\lambda
        \right) \subseteq \bigcap_{\lambda \in \Lambda}\vV(I_\lambda).
    \]
    Από την άλλη, τα στοιχεία του $\sum_{\lambda \in \Lambda}I_\lambda$ είναι της μορφής $g = \sum_{\lambda \in \Lambda} f_\lambda$, για κάποια $f_\lambda \in I_\lambda$. Επομένως, αν $\bfa \in \vV(I_\lambda)$, για κάθε $\lambda \in \Lambda$, τότε 
    \[
    g(\bfa) = \sum_{\lambda \in \Lambda} f_\lambda(\bfa) = 0
    \]
    και γι αυτό 
    \[
    \bigcap_{\lambda \in \Lambda}\vV(I_\lambda) \subseteq \vV\left(
            \sum_{\lambda \in \Lambda} I_\lambda
        \right)
    \]
    και το ζητούμενο έπεται.
\end{proof}

\begin{example}
    \label{ex:zero_locus_properties}
    Ας δούμε τι λέει η Πρόταση \ref{prop:zero_locus} με ένα παράδειγμα. Έστω $I = (y-x^2)$ και $J = (y-1)$ δύο ιδεώδη του $\RR[x,y]$. Τότε 
    \begin{align*}
    IJ &= \left((y-x^2)(y-1)\right) \\
    I+J &= (y-1, x^2-1),
    \end{align*}
    όπου η δεύτερη ισότητα έπεται από την παρατήρηση $y-x^2 = (y-1) - (x^2-1)$. Συνεπώς, 
    \[
    \begin{tikzpicture}
        \begin{axis}[
            axis lines = middle,
            axis line style={-latex, gray!50},
            xlabel = {$x$},
            ylabel = {$y$},
            label style = {gray!50},
            xmin=-2, xmax=2,
            ymin=-2, ymax=4,
            samples=100,
            xticklabels=\empty,
            yticklabels=\empty,
            ticks=none
        ]
        
            % function
            \addplot[burntorange, thick] {x^2}; 
            \addplot[iceberg, thick] {1}; 
        \end{axis}

        \begin{axis}[
            at={(8cm,0)},
            axis lines = middle,
            axis line style={-latex, gray!50},
            xlabel = {$x$},
            ylabel = {$y$},
            label style = {gray!50},
            xmin=-2, xmax=2,
            ymin=-2, ymax=4,
            samples=100,
            xticklabels=\empty,
            yticklabels=\empty,
            ticks=none
        ]
        
            % function
            \addplot[burntorange!20, thick] {x^2};  
            \addplot[iceberg!20, thick] {1}; 

            \addplot[
                only marks,
                mark=*,
                mark size=2pt
            ] coordinates {(-1,1) (1,1)};
        \end{axis}
    \end{tikzpicture}
    \]
    \begin{align*}
        \vV(I) \cup \vV(J) &= \vV\left((\mathcolor{burntorange}{y-x^2})(\mathcolor{iceberg}{y-1})\right) \\
        \vV(I) \cap \vV(J) &=  \vV(y-1, x^2-1) = \{(\pm1,1)\}.
    \end{align*}
\end{example}

\begin{remark}
Η Πρόταση \ref{prop:zero_locus} μαζί με τις ιδιότητες των Παραδειγμάτων \ref{ex:variety} (α) και (β) μας πληροφορούν ότι ο $\AA^n$ είναι τοπολογικός χώρος του οποίου τα \emph{κλειστά σύνολα} είναι ακριβώς οι πολλαπλότητες. Η τοπολογία αυτή ονομάζεται \emph{τοπολογία Zariski}. Σε αντίθεση με την συνήθη τοπολογία του $\RR$, όπου τα κλειστά υποσύνολα είναι τα κλειστά διαστήματα, το Παράδειγμα \ref{ex:variety} (γ) μας πληροφορεί ότι στην τοπολογία Zariski τα κλειστά υποσύνολα είναι τα $\emptyset, \AA_\RR^1$ και κάθε πεπερασμένη συλλογή στοιχείων του $F$. Στις σημειώσεις αυτές δε θα αρχοληθούμε με αυτή την πτυχή των πολλαπλοτήτων. 
% Για μια εισαγωγή στην γενική τοπολογία παραπέμπουμε στο \cite{Mun14}.
\end{remark}

Η απεικόνιση $I \mapsto \vV(I)$ παίρνει ένα ιδεώδες, το οποίο είναι αλγεβρικής φύσης, και το αντιστοιχεί σε ένα υποσύνολο του αφφινικού χώρου, το οποίο είναι σε γεωμετρικής φύσης. Μπορούμε να κάνουμε και το αντίστροφο.
\begin{definition}
    \label{def:vanishing_ideal}
    Έστω $A \subseteq \AA^n$. Το σύνολο\footnote{Στην βιβλιογραφία συμβολίζεται και ως $\mathbb{I}(\cdot)$.} 
    \[
    \iI(A) \coloneqq \{f \in \Pol{n} : f(\bfa) = 0, \ \text{για κάθε $\bfa \in A$}\}
    \]
    ονομάζεται \defn{ιδεώδες μηδενισμού} (vanishing ideal) του $A$.
\end{definition}

\begin{remark}
    \leavevmode
    \begin{itemize}
        \item Το ιδεώδες μηδενισμού ενός υποσυνόλου του $\AA^n$ είναι πράγματι ιδεώδες (γιατί;).
        \item Αν $A \subseteq B$, τότε $\iI(B) \subseteq \iI(A)$ (γιατί;). Με άλλα λόγια, όσο μεγαλύτερο είναι το $A$, τόσο μικρότερο είναι το ιδεώδες μηδενισμού του. Για παράδειγμα, αν $A = \{(0,0)\}$ και $B = \{(a,0) : a \in F\}$, τότε κάθε πολυώνυμο του $F[x,y]$  κάθε όρος του οποίου περιέχει το $y$ θα μηδενίζεται στο $B$ και γι αυτό $\iI(B) = (y)$. Από την άλλη, κάθε πολυώνυμο του $F[x,y]$ χωρίς σταθερό όρο μηδενίζεται στο $A$ και γι αυτό $\iI(A) = (x,y)$. Με άλλα λόγια,
        \[
        A = \{(0,0)\} \ \subset \ \{(a,0) : a \in F\} = B
        \]
        και
        \[
        \iI(A) = (x, y) \ \supset \ (y) = \iI(B).
        \]
    \end{itemize}
\end{remark}

Στο ύφος της δεύτερης παρατήρησης, ένα ακραίο παράδειγμα είναι 
\[
\iI(\emptyset) = \Pol{n}.
\]
(γιατί;) Στον αντίποδα, έχουμε το εξής.
\begin{example}
    \label{ex:vanishing_ideal_of_emptyset}
    Υπενθυμίζουμε ότι το $F$ είναι ένα άπειρο σώμα. Θα δείξουμε ότι 
    \[
    \iI(\AA_F^n) = 0.
    \]
    Με άλλα λόγια\footnote{Ταυτίζοντας τα πολυώνυμα στις $n$ μεταβλητές με τις πολυωνυμικές συναρτήσεις $\AA^n \to \AA^1$, όπως κάναμε στο τέλος της Ενότητας 3, η παρατήρηση αυτή μας πληροφορεί κάτι γεωμετρικό χρησιμοποιώντας άλγεβρα.}, κάθε μη μηδενικό πολυώνυμο σε $n$ μεταβλητές είναι μη μηδενικό σε κάποιο σημείο του $\AA^n$ ή ισοδύναμα αν ένα πολυώνυμο μηδενίζεται σε ολόκληρο το $\AA^n$, τότε πρέπει αναγκαστικά να είναι το μηδενικό πολυώνυμο.
\end{example}


% \begin{thebibliography}{99}
%     % \bibitem{DF04}
%     % D. S. Dummit και R. M. Foote, 
%     % \textit{Abstract algebra},
%     % third edition, Wiley, 2004. 
%     % \bibitem{Mun14}
%     % J.~Munkres, 
%     % \textit{Topology},
%     % second edition, Prentice Hall, 2000.
% \end{thebibliography}
\end{document}